\documentclass[../main.tex]{subfiles}
\begin{document}

\chapter{同期}
\lessoninfo{
  video={Lesson 20，動画 1--23},
  textbook={H\&P \cite{hennessy2017} §5.5，付録 I.4},
  keywords={ロック，相互排除，クリティカルセクション，スピンロック，Lamport のパン屋のアルゴリズム，アトミック命令，アトミック交換，テストアンドセット，LL/SC，リンクレジスタ，キャッシュコヒーレンス，バリア同期},
}

第19章では，すべてのコアのキャッシュが共有メモリについて1つのコヒーレント
な見え方を示すようにした．コヒーレントなメモリだけでは，共有メモリの
プログラムは正しくならない．複数のスレッドが同じ変数を読んで更新するとき，
結果はそれらの命令がどのように交錯するかに依存するからである．この章では，
そうした更新を一度に1つのスレッドだけに許すロックを導入し，通常のロード
とストアで作った素朴なロックがなぜ正しくないのか，プロセッサが代わりに
どのようなアトミック命令を提供するのかを示し，ロック変数がキャッシュ
コヒーレンスとどう相互作用するかを調べ，最後に，スレッドの一群を互いに
待ち合わせさせるバリアを扱う．

\section{同期の必要性}
\label{sec:l20-need}

メモリを共有するスレッドは，共有データの更新を協調させなければならない．
\source{Lesson 20，動画 1--2；H\&P §5.5．}例として，長い文書の中で各文字
が何回現れるかを数えるプログラムを考える．2つのスレッドが文書の半分ずつを
処理し，26 個のカウンタからなる共有配列をインクリメントする
（\cref{ex:l20-count}）．

\begin{example}[文字を数える]
  \label{ex:l20-count}
  各スレッドは，小文字だけからなる文書の自分の担当分に対して次のループを
  実行する．
  \begin{lstlisting}[language=C, numbers=none, morekeywords={shared}]
shared char doc[N];
shared int  count[26];
for (int i = lo; i < hi; i++) {  /* 自分の担当分 */
  int c = doc[i] - 'a';
  count[c]++;
}
\end{lstlisting}
\end{example}

文 \texttt{count[c]++} はソースコード 1 行だが，プロセッサにとっては
カウンタのロード，加算，ストアの 3 命令である．
\begin{lstlisting}[language={[x86masm]Assembler}, numbers=none]
LW   R1, 0(R2)      ; R2 = count[c] のアドレス
ADDI R1, R1, 1
SW   R1, 0(R2)
\end{lstlisting}
各スレッドは自分自身のレジスタを持つが，両方のスレッドが同じ文字に出会った
とき，それぞれの R2 は同じアドレスを保持している．どちらかがインクリメント
した値をストアする前に両方のスレッドがカウンタをロードすると，両者は同じ
新しい値を計算し，1回分のインクリメントが失われる（\cref{tab:l20-race}）．

\begin{table}[htbp]
  \centering\small
  \caption{2つのスレッドが，41 を保持する文字 \texttt{e} のカウンタを
    インクリメントする．2回のインクリメントの後，カウンタは 43 ではなく
    42 を保持している．}
  \label{tab:l20-race}
  \begin{tabular}{@{}cllrrr@{}}
    \toprule
    手順 & スレッド A & スレッド B & A の R1 & B の R1 & カウンタ \\
    \midrule
    1 & \texttt{LW R1, 0(R2)}   &                             & 41 & --  & 41 \\
    2 &                         & \texttt{LW R1, 0(R2)}       & 41 & 41  & 41 \\
    3 & \texttt{ADDI R1, R1, 1} &                             & 42 & 41  & 41 \\
    4 &                         & \texttt{ADDI R1, R1, 1}     & 42 & 42  & 41 \\
    5 & \texttt{SW R1, 0(R2)}   &                             & 42 & 42  & 42 \\
    6 &                         & \texttt{SW R1, 0(R2)}       & 42 & 42  & 42 \\
    \bottomrule
  \end{tabular}
\end{table}

インクリメントが失われるかどうかは，スレッド間の相対的なタイミングに依存
する．たいていは2つのスレッドが異なる文字を処理するか，一方のスレッドが
3命令を終えてからもう一方が始めるので，カウントは正しい．ときおりカウント
が小さく出るが，その誤りはプログラムを再度実行しても再現するとは限らない．
\margin{時分割された1つのコアでも，ロードとストアの間にコンテキスト
スイッチが入れば同じ損失が起こる．}したがってロード，加算，ストアは，同じ
カウンタを更新する他のあらゆるスレッドから見て1つの分割不可能な単位として
実行されなければならない．これを提供するのが\emph{同期}の役割である．

\section{ロック}
\label{sec:l20-locks}

同期の最も基本的な機構がロックである．\source{Lesson 20，動画
3--4；H\&P §5.5．}

\begin{definition}[ロック]
  \term[ろっく]{ロック}[lock]とは，2つの操作を持つ共有オブジェクトで
  ある．\texttt{lock} はロックが空くまで待ち，空いたら呼び出したスレッドが
  保持していると印を付ける．\texttt{unlock} は再び空きの印を付ける．任意の
  時点で，あるロックを保持するスレッドは高々1つである．
\end{definition}

カウンタごとに1つのロックを置くと，\cref{ex:l20-count} のループは次のよう
になる．
\begin{lstlisting}[language=C, numbers=none, morekeywords={shared}]
shared lock_t countLock[26];
  ...
  lock(&countLock[c]);
  count[c]++;
  unlock(&countLock[c]);
\end{lstlisting}
スレッド A がロックを保持している間にスレッド B が \texttt{lock} を呼ぶと，
B は A が \texttt{unlock} を呼ぶまで待ち，そのあとで初めてカウンタを
ロードする（\cref{fig:l20-lock-timing}）．したがってスレッド B はスレッド~A
がストアした値を見るので，インクリメントは失われない．

\begin{figure}[htbp]
  \centering
  % Two threads that use the same lock.  Thread B calls lock() while thread A
% is in its critical section, waits, and enters when A calls unlock().
\begin{tikzpicture}[x=0.9cm,y=1cm, font=\footnotesize\sffamily, >={Stealth[length=1.6mm]},
  ev/.style={font=\scriptsize\ttfamily, inner sep=1pt}]
  \def\lnbox#1#2#3#4{% y, x0, x1, fill
    \draw[fill=#4] (#2,{#1-0.25}) rectangle (#3,{#1+0.25});}
  % ---- thread A
  \node[anchor=east] at (-0.15,1) {\iflangja{スレッド A}{thread A}};
  \lnbox{1}{0}{1.5}{white}
  \lnbox{1}{1.5}{4.5}{ln-accent!20}
  \lnbox{1}{4.5}{10}{white}
  \node[font=\scriptsize] at (3,1) {\iflangja{クリティカルセクション}{critical section}};
  \node[ev, anchor=south] at (1.5,1.3) {lock()};
  \node[ev, anchor=south] at (4.5,1.3) {unlock()};
  \draw (1.5,1.25) -- (1.5,1.32);
  \draw (4.5,1.25) -- (4.5,1.32);
  % ---- thread B
  \node[anchor=east] at (-0.15,0) {\iflangja{スレッド B}{thread B}};
  \lnbox{0}{0}{2.5}{white}
  \lnbox{0}{2.5}{4.5}{ln-caution!15}
  \lnbox{0}{4.5}{7.5}{ln-accent!20}
  \lnbox{0}{7.5}{10}{white}
  \node[font=\scriptsize] at (3.5,0) {\iflangja{待機}{waiting}};
  \node[font=\scriptsize] at (6,0) {\iflangja{クリティカルセクション}{critical section}};
  \node[ev, anchor=north] at (2.5,-0.3) {lock()};
  \node[ev, anchor=north] at (7.5,-0.3) {unlock()};
  \draw (2.5,-0.25) -- (2.5,-0.32);
  \draw (7.5,-0.25) -- (7.5,-0.32);
  % ---- hand-over of the lock
  \draw[->, ln-caution] (4.5,0.75) -- (4.5,0.25);
  % ---- time axis
  \draw[->, ln-muted] (0,-0.95) -- (10.3,-0.95)
    node[right, inner sep=2pt] {\iflangja{時間}{time}};
\end{tikzpicture}

  \caption{同じロックを使う2つのスレッド．スレッド B は，スレッド A が
    クリティカルセクションを出るまで \texttt{lock} の中で待つ．}
  \label{fig:l20-lock-timing}
\end{figure}

\begin{definition}[相互排除]
  あるロックで保護されたコードを一度に高々1つのスレッドしか実行しないと
  いう保証を\term[そうごはいじょ]{相互排除}[mutual exclusion]という．
\end{definition}

このためロックは\emph{ミューテックス}%
\index{みゅーてっくす@ミューテックス|see{ロック}}とも呼ばれる．

\begin{definition}[クリティカルセクション]
  \term[くりてぃかるせくしょん]{クリティカルセクション}[critical section]
  とは，相互排除のもとで実行されなければならない命令の並びであり，
  \texttt{lock} で始まり \texttt{unlock} で終わる．
\end{definition}

ロックが保護するのはコードではなく\emph{データ}である．相互排除が成り立つ
のは同じロックを使うスレッドの間だけなので，ある共有変数へのアクセスは
すべて同一のロックのもとで行わなければならない．異なるロックを保持したまま
同じ変数を更新するスレッドどうしは，まったく排他しない．1つのロックが
いくつの変数を保護するかは，性能に影響する選択である．\tip{ある共有変数
へのアクセスにはすべて同じロックを使い，独立な変数には異なるロックを使う
こと．}26 個のカウンタすべてに1つのロックを使っても正しいが，その場合は
どのインクリメントも，別の文字のものであってさえ他のすべてを排他するので，
スレッドは時間の多くを互いの待ち合わせに費やす．カウンタごとに1つのロックを
置けば，異なるカウンタをインクリメントするスレッドは並列に進み，あるスレッド
が待つのは別のスレッドが同じカウンタを更新している間だけである．

\section{ロックの実装}
\label{sec:l20-implementing}

ロックは，空いているときは 0，保持されているときは 1 をとる共有の
\emph{ロック変数}で表せる．\source{Lesson 20，動画 5--7；H\&P §5.5；
\cite{lamport1974}．}素朴な実装は，この変数が 0 になるまで待ち，そのあと
それを~1 にする．
\begin{lstlisting}[language=C, numbers=none]
void lock(lock_t *l) {
  while (*l == 1)
    ;              /* ロックが空くまで待つ */
  *l = 1;          /* 取得する */
}
void unlock(lock_t *l) {
  *l = 0;
}
\end{lstlisting}
待っているスレッドはこのループを実行し続け，ロック変数を何度も読む．この
ように待つロックを\term[すぴんろっく]{スピンロック}[spin lock]という．

\needspace{7\baselineskip}
この実装は動作しない．機械語では，ロック変数を検査することと，それに値を
設定することは別々の命令だからである．
\begin{lstlisting}[language={[x86masm]Assembler}, numbers=none, deletekeywords={lock}]
lock:  LW   R1, 0(R2)    ; R2 = ロック変数のアドレス
       BNEZ R1, lock     ; 保持されている: もう一度読む
       ADDI R1, R0, 1
       SW   R1, 0(R2)    ; ロック変数 <- 1
\end{lstlisting}
ロックが空いている間に2つのスレッドがともに \texttt{LW} を実行すると，
両者とも 0 を読み，ともにループを抜け，ともに 1 をストアし，ともに
クリティカルセクションに入ってしまう．これは \cref{sec:l20-need} で見た更新の
消失にほかならない．検査と設定は共有変数の読み出しと更新であり，したがって
\texttt{lock} 自身がクリティカルセクションを必要とすることになる---それこそ
\texttt{lock} が提供すべきものである．

逃れ道は2つある．
\begin{itemize}
  \item \textbf{ソフトウェア．}ソフトウェアによる解法として
    \term[らんぽーとのぱんやのあるごりずむ]{Lamport のパン屋の
    アルゴリズム}[Lamport's bakery algorithm]\cite{lamport1974} があり，
    通常のロードとストアだけで相互排除を実現する．これは複雑であり，1回の
    取得ごとに多数の命令を実行するためロックが遅くなる---きわめて頻繁に
    取得されるロックには遅すぎる．
  \item \textbf{ハードウェア．}プロセッサが，ロック変数の読み出しと
    書き込みを1つの分割不可能な手順で行う命令を提供する．
\end{itemize}

\begin{definition}[アトミック命令]
  \term[あとみっくめいれい]{アトミック命令}[atomic instruction]とは，
  複数の手順---ここではメモリ位置の読み出しとその書き込み---を1つの分割
  不可能な操作として行う命令である．このコアであれ他のどのコアであれ，
  手順の間にその位置へアクセスできる命令は存在しない．
\end{definition}

ロックを実装するには，アトミックな手順がロック変数の読み出しと書き込みの
両方を含まなければならない．読み出した値の検査は，そのあと通常の命令で
行ってよい．いったん変数が書かれてしまえば，他のどのスレッドも古い値を
読むことはできないからである．

\section{アトミックな読み書き命令}
\label{sec:l20-atomic}

命令セットは，メモリ位置の読み出しと書き込みを1つのアトミックな手順で行う
命令を複数提供している．\source{Lesson 20，動画 8--10；H\&P §5.5．}その中で
最も単純なものは，レジスタとメモリを入れ替える．

\begin{definition}[アトミック交換]
  \term[あとみっくこうかん]{アトミック交換}[atomic exchange]
  \texttt{EXCH R1, 0(R2)} は，R1 の内容を R2 のアドレスにあるメモリ位置へ
  アトミックに書き込み，その位置の元の内容を R1 に残す．
\end{definition}

\needspace{6\baselineskip}
スレッドは，0 が返るまでロック変数に 1 を交換で書き込むことによって，
ロックを取得する．
\begin{lstlisting}[language={[x86masm]Assembler}, numbers=none, deletekeywords={lock}]
        ADDI R1, R0, 1
lock:   EXCH R1, 0(R2)   ; R1 とロック変数を交換
        BNEZ R1, lock    ; 保持されていた: やり直す
        ...              ; クリティカルセクション
unlock: SW   R0, 0(R2)   ; ロック変数 <- 0
\end{lstlisting}
ロックが保持されていれば，交換は 1 の上に 1 を書いて 1 を返すので，R1 は
次の試行に向けて再び 1 になる．ロックが空いていれば，最初に実行された交換が
0 を返して 1 を残し，それ以降の交換はすべて 1 を返す．読み出しと書き込みの
間には何も割り込めないので，2つのスレッドがともに~0 を読むことはありえない．

交換は試行のたびにロック変数に書き込む．他のスレッドがロックを保持している
間もそうであり，そこでは書き込みは何も変えない．この代償については
\cref{sec:l20-performance} で示す．2つ目の命令は，読み出した値が検査を
通ったときだけ書き込む．

\begin{definition}[テストアンドセット]
  \term[てすとあんどせっと]{テストアンドセット}[test-and-set]命令は，
  メモリ位置を検査し，その値が検査を通った場合にのみそこへ書き込む．ここで
  の検査は 0 との比較である．\texttt{TSTSW R1, 0(R2)} は，
  $\text{Mem}[\text{R2}] = 0$ ならば
  $\text{Mem}[\text{R2}] \leftarrow \text{R1}$ と
  $\text{R1} \leftarrow 1$ とし，そうでなければメモリを変えずに
  $\text{R1} \leftarrow 0$ とする．
\end{definition}

\needspace{7\baselineskip}
この命令は書き込みを行ったかどうかを R1 で報告するので，R1 が~1 で返れば
ロックを取得できている．試行に失敗したスレッドは，再試行の前にストアすべき
値を R1 に入れ直さなければならない．
\begin{lstlisting}[language={[x86masm]Assembler}, numbers=none, deletekeywords={lock}]
lock:   ADDI R1, R0, 1   ; ロック変数へストアする値
        TSTSW R1, 0(R2)  ; ロック変数が 0 ならストア
        BEQZ R1, lock    ; ストアしなかった: やり直す
        ...              ; クリティカルセクション
unlock: SW   R0, 0(R2)   ; ロック変数 <- 0
\end{lstlisting}
他のスレッドがロックを保持している間は検査が失敗し，命令は何も書き込まない．

どちらの命令も，ロード/ストア型の命令セットには収まりが悪い．他のすべての
命令はデータメモリに高々1回しかアクセスしない．ロードは読み，ストアは書き，
いずれもパイプラインの単一の MEM ステージで行われる（第3章）．交換や
テストアンドセットは，1つの命令の中で同じ位置を読み\emph{かつ}書くので，
パイプラインにはメモリアクセスを2回行うデータパス---事実上2つ目のメモリ
ステージ---が必要になり，さらにプロセッサは，他のコアによるアクセスも例外
も，その間に入らないことを保証しなければならない．

\section{LL/SC}
\label{sec:l20-llsc}

もう1つの方法は，アトミックな読み出しと書き込みを2つの命令に分け，その
それぞれが通常のロードやストアと同じくメモリに1回だけアクセスするように
することである．\source{Lesson 20，動画 11--13；H\&P §5.5．}

\begin{definition}[LL/SC]
  命令対 \term{LL/SC}[load-linked/store-conditional] は，次の2つの命令から
  なる．
  \begin{description}
    \item[\texttt{LL R1, 0(R2)}] は \texttt{LW} と同じく
      $\text{R1} \leftarrow \text{Mem}[\text{R2}]$ とロードし，その
      アドレスをコアの\term[りんくれじすた]{リンクレジスタ}[link register]
      に記録する．
    \item[\texttt{SC R1, 0(R2)}] は \texttt{SW} と同じく
      $\text{Mem}[\text{R2}] \leftarrow \text{R1}$ とストアするが，
      リンクレジスタが依然としてこのアドレスを保持している場合に限る．
      その場合は成功を報告するために R1 を 1 にする．そうでなければ何も
      ストアせず，R1 を~0 にする．
  \end{description}
\end{definition}

この対をアトミックにしているのは，リンクレジスタを消去する規則である．
メモリ位置への書き込みは---\texttt{SC} であれ通常のストアであれ，この
コアによるものであれ他のどのコアによるものであれ---その位置のアドレスを
保持しているすべてのリンクレジスタを消去する（\cref{fig:l20-llsc}）．
\margin{H\&P によれば，\texttt{LL} と \texttt{SC} の間の割り込みや例外も
リンクレジスタを消去する．}あるコアは，第19章のコヒーレンスプロトコルを
通じて他のコアの書き込みを知る．リンクされたアドレスを含むブロックの
無効化が，リンクレジスタを消去するのである．したがって検出はブロック単位で
働き，同じブロック内の別のアドレスへの書き込みもリンクを消去する．これは
再試行の代償を払うだけで，正しさを損なうことはない．

\texttt{SC} が成功するのは，対応する \texttt{LL} 以降に誰もその位置に
書き込んでいない場合だけである．とくにそのとき，\texttt{LL} が読んだ値は
まだメモリにある．逆は成り立たない．同じ値をもう一度書いても \texttt{SC} は
失敗する．LL/SC が検出するのは書き込みであって値の変化ではないからである．
\texttt{LL} と \texttt{SC} の間には他の命令が走ってよく，同じ位置に
アクセスする他のスレッドの命令も含まれる．それらが干渉すれば \texttt{SC}
がそれを検出して失敗し，プログラムは単にやり直す．どちらの命令も単独では
通常のロードやストア以上のことをしないにもかかわらず，効果はアトミックな
読み出しと書き込みのそれである．

\needspace{9\baselineskip}
LL/SC で作ったロックは，ロック変数をロードして検査し，それが~0 だった場合に
のみ 1 のストアを試みる．\tip{\texttt{LL} と \texttt{SC} の間に置いてよいのは
レジスタ間の命令だけである．別のメモリアクセスはリンクを消去しうる．}
\begin{lstlisting}[language={[x86masm]Assembler}, numbers=none, deletekeywords={lock}]
lock:   LL   R1, 0(R2)   ; R1 <- ロック変数，リンクを記録
        BNEZ R1, lock    ; 保持されている: やり直す
        ADDI R1, R0, 1   ; R1 <- 1
        SC   R1, 0(R2)   ; リンクが有効なら 1 をストア
        BEQZ R1, lock    ; SC 失敗: 最初からやり直す
        ...              ; クリティカルセクション
unlock: SW   R0, 0(R2)   ; ロック変数 <- 0
\end{lstlisting}
\begin{figure}[htbp]
  \centering
  % LL/SC on two cores.  Both cores load-link the lock variable at address A;
% core 0's store-conditional succeeds, and its store clears every link
% register that holds A, so core 1's store-conditional fails.
% Register contents are shown after step (5).
\begin{tikzpicture}[x=1cm,y=1cm, font=\footnotesize\sffamily, >={Stealth[length=1.8mm]},
  core/.style={draw, fill=ln-accent!15, minimum width=3.6cm, minimum height=0.55cm, inner sep=2pt},
  reg/.style={draw, fill=white, minimum width=1.8cm, minimum height=0.5cm, inner sep=1pt, font=\footnotesize\ttfamily},
  hd/.style={inner sep=1pt, font=\scriptsize\sffamily, text=ln-muted},
  big/.style={draw, fill=ln-box, minimum height=0.65cm, inner sep=3pt},
  op/.style={font=\scriptsize, align=left, anchor=south west, inner sep=1pt}]
  % ---- core 0
  \node[core] (c0) at (1.8,0) {\iflangja{コア}{core}~0};
  \node[reg, anchor=north west] (r0) at (c0.south west) {1};
  \node[reg, anchor=north east] (l0) at (c0.south east) {--};
  \node[hd, anchor=north] (r0l) at (r0.south) {R1};
  \node[hd, anchor=north] (l0l) at (l0.south) {\iflangja{リンクレジスタ}{link register}};
  % ---- core 1
  \node[core] (c1) at (8.4,0) {\iflangja{コア}{core}~1};
  \node[reg, anchor=north west] (r1) at (c1.south west) {0};
  \node[reg, anchor=north east, fill=ln-caution!15] (l1) at (c1.south east) {--};
  \node[hd, anchor=north] (r1l) at (r1.south) {R1};
  \node[hd, anchor=north] (l1l) at (l1.south) {\iflangja{リンクレジスタ}{link register}};
  % ---- operations, numbered in the order they happen
  \node[op] at ([yshift=2pt]c0.north west) {%
    \iflangja{(1) LL: R1 $=$ 0，リンクレジスタ $\leftarrow$ A}%
             {(1) LL: R1 $=$ 0, link register $\leftarrow$ A}\\
    \iflangja{(3) SC: リンク有効 $\to$ 1 を格納，R1 $=$ 1}%
             {(3) SC: link valid $\to$ store 1, R1 $=$ 1}};
  \node[op] at ([yshift=2pt]c1.north west) {%
    \iflangja{(2) LL: R1 $=$ 0，リンクレジスタ $\leftarrow$ A}%
             {(2) LL: R1 $=$ 0, link register $\leftarrow$ A}\\
    \textcolor{ln-caution}{\iflangja{(5) SC: リンク無効 $\to$ 格納せず，R1 $=$ 0}%
             {(5) SC: link cleared $\to$ no store, R1 $=$ 0}}};
  % ---- memory
  \node[big, minimum width=4.4cm] (mem) at (5.1,-2.55)
    {\iflangja{ロック変数（アドレス A）$=$ 1}{lock variable (address A) $=$ 1}};
  % store of core 0
  \draw[->] (r0l.south) |- (mem.west)
    node[pos=0.75, above, font=\scriptsize] {\iflangja{(3) 1 を格納}{(3) store 1}};
  % clearing of the link registers
  \draw[->, ln-caution] (mem.east) -- (11.0,-2.55) |- (l1.east)
    node[pos=0, above left, font=\scriptsize, align=right, inner sep=2pt]
    {\iflangja{(4) A への書き込みで，\\A を保持する\\リンクレジスタを消去}%
              {(4) the write to A clears\\every link register holding A}};
\end{tikzpicture}

  \caption{2つのコアでの LL/SC．コア 0 のストアが成功すると，コア 1 の
    リンクレジスタが消去され，コア 1 の SC は失敗する．}
  \label{fig:l20-llsc}
\end{figure}


\needspace{6\baselineskip}
\begin{example}[1つのロックを奪い合う2つのコア]
  \label{ex:l20-llsc}
  コア 0 とコア 1 が上の \texttt{lock} のコードを同時に実行する．アドレス A
  にあるロック変数は 0 である．\cref{tab:l20-llsc} は，命令を実行される順に
  追ったものである．どちらのコアもロックが空いていると判断するが，成功
  できる \texttt{SC} は一方だけである．そのストアが他方のコアのリンク
  レジスタを消去するからである．コア 1 は \texttt{LL} に戻り，今度は 1 を
  読んで待つ．
\end{example}

\begin{table}[htbp]
  \centering\small
  \caption{\cref{ex:l20-llsc} のトレース．リンクレジスタの列のダッシュは，
    そのレジスタがアドレスを保持していないことを表す．}
  \label{tab:l20-llsc}
  \begin{tabular}{@{}ccllccc@{}}
    \toprule
    & & & & \multicolumn{2}{c}{リンクレジスタ} & ロック \\
    \cmidrule(lr){5-6}
    手順 & コア & 命令 & 効果 & コア 0 & コア 1 & 変数 \\
    \midrule
    1  & 0 & \texttt{LL R1, 0(R2)}  & R1 $=$ 0                  & A  & -- & 0 \\
    2  & 1 & \texttt{LL R1, 0(R2)}  & R1 $=$ 0                  & A  & A  & 0 \\
    3  & 0 & \texttt{BNEZ}, \texttt{ADDI} & R1 $=$ 1            & A  & A  & 0 \\
    4  & 1 & \texttt{BNEZ}, \texttt{ADDI} & R1 $=$ 1            & A  & A  & 0 \\
    5  & 0 & \texttt{SC R1, 0(R2)}  & 1 をストア，R1 $=$ 1      & -- & -- & 1 \\
    6  & 1 & \texttt{SC R1, 0(R2)}  & ストアなし，R1 $=$ 0      & -- & -- & 1 \\
    7  & 0 & \texttt{BEQZ R1, lock} & 不成立: ロックを保持      & -- & -- & 1 \\
    8  & 1 & \texttt{BEQZ R1, lock} & 成立: 最初からやり直す    & -- & -- & 1 \\
    9  & 1 & \texttt{LL R1, 0(R2)}  & R1 $=$ 1                  & -- & A  & 1 \\
    10 & 1 & \texttt{BNEZ R1, lock} & 成立: ロックは保持済み    & -- & A  & 1 \\
    \bottomrule
  \end{tabular}
\end{table}

\needspace{9\baselineskip}
1回の読み出し-変更-書き込みだけからなるクリティカルセクションは，LL/SC が
使えるならロックをまったく必要としない．対を共有変数そのものに適用すれば
よい．\cref{ex:l20-count} のカウンタであれば，
\begin{lstlisting}[language={[x86masm]Assembler}, numbers=none]
retry:  LL   R1, 0(R2)   ; R2 = count[c] のアドレス
        ADDI R1, R1, 1
        SC   R1, 0(R2)   ; 誰も書いていなければストア
        BEQZ R1, retry   ; 誰かが書いた: 読み直して加算
\end{lstlisting}
がカウンタをアトミックにインクリメントする．\texttt{LL} と \texttt{SC} の
間に別のスレッドがカウンタに書き込めば \texttt{SC} が失敗し，新しい値で
インクリメントがやり直されるので，更新は失われない．

\begin{keypoint}
  ロックには，ロック変数の読み出しと書き込みが1つのアトミックな手順を
  なすことが必要である．アトミック交換とテストアンドセットは，その両方を
  1つの変則的な命令で行う．LL/SC は通常のメモリアクセス2回を用い，その間に
  誰かがその位置に書いていればストアを失敗させる．これは，ロックなしで短い
  読み出し-変更-書き込みをアトミックにもする．
\end{keypoint}

\section{ロックと性能}
\label{sec:l20-performance}

ロック変数も他の共有変数と変わらない．ロックを使うどのコアも，それを含む
ブロックを自分のキャッシュに持ち，第19章のキャッシュコヒーレンス
プロトコルがこれらのコピーをコヒーレントに保つ．\source{Lesson 20，動画
14--17；H\&P §5.5．}アトミック交換の上に作ったスピンロックは，コヒーレンス
との相性が悪い．

\texttt{EXCH} は，ロックが保持されているときでもロック変数に書き込む．1 の
上に 1 を書いても値は変わらないが，コヒーレンスプロトコルにとっては他と
変わらない書き込みである．MSI プロトコルのもとでは，書き込むコアはその
ブロックを M（modified）状態で持たなければならず，他のすべてのコピーは無効化
される．1つのコアがロックを保持している間，\texttt{EXCH} を実行する待機中の
各コアは，その直前に試行したコアからブロックを奪う．こうしてほとんどすべて
の試行がコヒーレンスミスを引き起こす．このコアの前回の試行以降に別の待機中
のコアがそのブロックに書き込んでいるからであり，ブロックは再びバス上を
運ばれる（\cref{tab:l20-coherence} 上段）．待機中のコアは何も成し遂げない
にもかかわらず，そのバストラフィックはエネルギーを消費し，他のすべてのコア
のメモリアクセスを遅くする．クリティカルセクションを実行しているロック
保持者もその1つであり，しかもロックを解放するときにブロックを取り戻さ
なければならない．

\needspace{7\baselineskip}
対策は，まず通常のロードでロック変数を検査し，ロックが空いているように
見えるときにだけアトミック命令を使うことである．
\begin{lstlisting}[language={[x86masm]Assembler}, numbers=none, deletekeywords={lock}]
lock:   LW   R1, 0(R2)   ; ロック変数の通常の読み出し
        BNEZ R1, lock    ; 保持されている: 読み続ける
        ADDI R1, R0, 1
        EXCH R1, 0(R2)   ; 空きに見える: 取得を試みる
        BNEZ R1, lock    ; 他のスレッドが取った: やり直す
\end{lstlisting}
これで待機中のコアはロック変数を読むだけになるので，どのコアもブロックを
S（shared）状態で保持し続け，そのロードはキャッシュヒットになる．ロックが保持
されている間，待機中のコアはバストラフィックをまったく発生させない
（\cref{tab:l20-coherence} 下段）．書き込みと，それに伴う無効化が起こるのは
解放の前後だけである．すなわち \texttt{unlock} のストアと，ロックが空いて
いると判断したコアによるアトミックな試行であり，成功するのはそのうち1つ
だけである．

\begin{table}[htbp]
  \centering\small
  \caption{コア 0 がロックを保持している間の，4つのコアのキャッシュに
    おけるロック変数を含むブロックの MSI 状態（M，S，I: modified，shared，
    invalid）．上段: コア 1--3 が \texttt{EXCH} でスピンする場合．下段:
    まず \texttt{LW} で検査する場合．}
  \label{tab:l20-coherence}
  \setlength{\tabcolsep}{4pt}
  \begin{tabular}{@{}lccccl@{}}
    \toprule
    アクセス & C0 & C1 & C2 & C3 & バストラフィック \\
    \midrule
    \multicolumn{6}{@{}l}{\emph{アトミック交換でスピンする場合}} \\
    初め: C0 がロックを保持                     & M & I & I & I & \\
    C1: \texttt{EXCH}（1 を読む）               & I & M & I & I & ミス，ブロックが移動 \\
    C2: \texttt{EXCH}（1 を読む）               & I & I & M & I & ミス，ブロックが移動 \\
    C3: \texttt{EXCH}（1 を読む）               & I & I & I & M & ミス，ブロックが移動 \\
    C1: \texttt{EXCH}（1 を読む）               & I & M & I & I & ミス，ブロックが移動 \\
    \quad\dots\ 試行ごとに1回のミスが続き，      &   &   &   &   & \\
    C0: \texttt{SW R0}（unlock）                & M & I & I & I & ミス，ブロックが移動 \\
    \midrule
    \multicolumn{6}{@{}l}{\emph{まず検査する場合}} \\
    初め: C0 がロックを保持                     & M & I & I & I & \\
    C1: \texttt{LW}（1 を読む）                 & S & S & I & I & ミス \\
    C2: \texttt{LW}（1 を読む）                 & S & S & S & I & ミス \\
    C3: \texttt{LW}（1 を読む）                 & S & S & S & S & ミス \\
    C1--C3: \texttt{LW} を繰り返す              & S & S & S & S & なし（ヒット） \\
    C0: \texttt{SW R0}（unlock）                & M & I & I & I & 無効化 \\
    C2: \texttt{LW}（0 を読む）                 & S & I & S & I & ミス \\
    C2: \texttt{EXCH}（0 を読んで取得）         & I & I & M & I & 無効化 \\
    \bottomrule
  \end{tabular}
\end{table}

先に検査するロックについて，誤りやすい点が2つある．
\begin{itemize}
  \item \textbf{ロードはフィルタにすぎない．}相互排除は依然としてアトミック
    命令によってもたらされる．ロックが空いていると判断した \texttt{LW} と
    \texttt{EXCH} の間に別のスレッドがロックを取得しうるので，\texttt{EXCH}
    が返す値はやはり検査しなければならない．\texttt{LW} の結果だけで
    クリティカルセクションへ分岐するロックは，\cref{sec:l20-implementing}
    の壊れたロックそのものである．
  \item \textbf{他の命令はそれ自体が先に検査している．}\texttt{LL} は
    読み出しであり，\cref{sec:l20-llsc} のロックは \texttt{LL} がロックの
    空きを確認したあとでのみ \texttt{SC} を試みる．テストアンドセットは，
    ロックが空いていると分かったときだけ書き込む．1 を読み続ける待機中の
    コアは決して書き込まないので，余分なロードは要らない．
\end{itemize}

最後に \texttt{unlock} はアトミック命令を必要としない．0 を通常のストアで
書けば十分である．\texttt{unlock} を実行するのはロックを保持しているスレッド
だけなので，2つのスレッドが同じロックを同時に解放することはない．その瞬間に
ロックを取得しようとしているスレッドは，ストアの前なら 1 を読んで待ち続け，
ストアの後なら 0 を読んでロックを奪い合う---どちらも正しい結果である．
LL/SC の場合，このストアは，直前にその変数を \texttt{LL} したコアのリンク
レジスタも消去する．これらのコアは 1 を読んでいて \texttt{SC} を試みないので，
影響を受けない．

\needspace{9\baselineskip}
\begin{keypoint}
  正しさは，ロック変数の検査と書き込みが1つのアトミックな手順であることを
  要求し，性能は書き込みが稀であることを要求する．待機中のスレッドはロック
  変数を読むだけにすべきであり，そうすればそのブロックはすべてのキャッシュ
  で共有状態に留まる．
\end{keypoint}

\section{バリア同期}
\label{sec:l20-barrier}

ロックはスレッドに順番を取らせる．\source{Lesson 20，動画 18--22；H\&P
付録 I.4．}これとは別の形の協調は，スレッドを互いに待ち合わせさせる．

\begin{definition}[バリア同期]
  \term[ばりあどうき]{バリア同期}[barrier synchronization]では，$N$ 個の
  スレッドからなる一群が，プログラム中の共通の地点で\emph{バリア}を呼ぶ．
  バリアに到達したスレッドはそこで待ち，$N$ 個のスレッドのうち最後のものが
  到着した時点で，全員が先へ進む（\cref{fig:l20-barrier}）．
\end{definition}

\begin{figure}[htbp]
  \centering
  % Barrier synchronization with three threads over two rounds.  A thread that
% reaches the barrier waits until the last thread of its round arrives; then
% all three continue together.
\begin{tikzpicture}[x=0.85cm,y=1cm, font=\footnotesize\sffamily, >={Stealth[length=1.6mm]}]
  \def\lnbox#1#2#3#4{% y, x0, x1, fill
    \draw[fill=#4] (#2,{#1-0.2}) rectangle (#3,{#1+0.2});}
  % releases of the barrier
  \foreach \x/\r in {5/{0 $\to$ 1}, 9.5/{1 $\to$ 0}} {
    \draw[dashed, ln-accent, thick] (\x,-0.55) -- (\x,2.45);
    \node[font=\scriptsize, text=ln-accent, anchor=south, align=center] at (\x,2.45)
      {\iflangja{解放}{released}\\\texttt{release}: \r};
  }
  % thread/y/arrival in round 0/arrival in round 1
  \foreach \t/\y/\a/\b in {0/2/2.5/9.5, 1/1/4/7.5, 2/0/5/8} {
    \node[anchor=east] at (-0.15,\y) {\iflangja{スレッド}{thread}~\t};
    \lnbox{\y}{0}{\a}{white}
    \ifdim\a pt<5pt \lnbox{\y}{\a}{5}{ln-caution!15}\fi
    \lnbox{\y}{5}{\b}{white}
    \ifdim\b pt<9.5pt \lnbox{\y}{\b}{9.5}{ln-caution!15}\fi
    \lnbox{\y}{9.5}{10.5}{white}
    \fill (\a,{\y+0.2}) -- ++(-0.09,0.16) -- ++(0.18,0) -- cycle;
    \fill (\b,{\y+0.2}) -- ++(-0.09,0.16) -- ++(0.18,0) -- cycle;
  }
  % time axis
  \draw[->, ln-muted] (0,-0.55) -- (10.8,-0.55)
    node[right, inner sep=2pt] {\iflangja{時間}{time}};
  % legend
  \begin{scope}[yshift=-1.0cm]
    \fill (0.09,-0.02) -- ++(-0.09,0.16) -- ++(0.18,0) -- cycle;
    \node[anchor=west, font=\scriptsize] at (0.25,0.05) {\iflangja{バリアに到着}{reaches the barrier}};
    \draw[fill=white] (3.2,-0.07) rectangle ++(0.35,0.25);
    \node[anchor=west, font=\scriptsize] at (3.6,0.05) {\iflangja{計算}{computing}};
    \draw[fill=ln-caution!15] (5.4,-0.07) rectangle ++(0.35,0.25);
    \node[anchor=west, font=\scriptsize] at (5.8,0.05) {\iflangja{バリアで待機}{waiting at the barrier}};
  \end{scope}
\end{tikzpicture}

  \caption{3つのスレッドが連続する2回のラウンドで使うバリア．各ラウンドは，
    そのラウンドの最後のスレッドが到着した時点で解放される．}
  \label{fig:l20-barrier}
\end{figure}

バリアは，段階を追って進む並列プログラムに適する．各スレッドが1つの段階の
自分の担当部分を計算し，次の段階が全スレッドの結果を必要とする場合である．
第18章の共有メモリによる総和はその一例である．スレッド 0 が合計を印字して
よいのは，すべてのスレッドが自身の部分和を加え終えたあと，すなわちバリアの
あとだけである．

\subsection{単純な実装}

バリアは2つの共有変数を使う．到着したスレッドの数を数えるカウンタと，待機中
のスレッドがスピンする\emph{解放フラグ}である．
\begin{lstlisting}[language=C, numbers=none, morekeywords={shared}]
shared int    count = 0, release = 0;
shared lock_t countLock;
void barrier(void) {
  lock(&countLock);
  if (count == 0) release = 0;  /* 最初: 閉じる */
  int arrived = ++count;        /* 自分を数える */
  unlock(&countLock);
  if (arrived == N) {           /* 最後: 開く */
    count = 0;
    release = 1;
  } else
    while (release == 0)
      ;                         /* 待つ */
}
\end{lstlisting}
カウンタにはロックが必要である．インクリメントが失われれば
（\cref{tab:l20-race}），カウンタは永遠に $N$ に届かない．ラウンドの最初の
スレッドが解放フラグをクリアするので，このラウンドのスレッドは待つことに
なる．最後のスレッドは次のラウンドに備えてカウンタをリセットし，フラグを
セットする．これが他のすべてのスレッドの待機ループを終わらせる．

\subsection{単純な実装が失敗する理由}

バリアは通常，ループの各段階で呼ばれるが，このように再利用すると上の実装は
デッドロックする．フラグは，他のスレッドがまだそれをスピンして読んでいる
最中にセットされるのであり，フラグが再び変わる前にそれを読むことを強制する
ものは何もない．あるスレッドが割り込みの処理中で，解放に気付くのが遅れたと
しよう．他のスレッドはバリアを出て，プログラムの次の段階を計算し，再び
バリアを呼ぶ．そのうち最初に到着したスレッドはカウンタが 0 であるのを見て，
解放フラグをクリアする．すると遅れたスレッドは 0 を読み，とうに解放された
バリアで待ち続ける．一方，他の $N-1$ 個のスレッドは，決して到着しない
スレッドを次のバリアで待つ．$N$ 個のスレッドすべてが永遠に待つことになる．

したがって待機中のスレッドは，フラグを見るのがどれほど遅れても，
\emph{自分の}ラウンドが解放されたことを認識できなければならない．

\subsection{再利用できるバリア}

再利用できるバリアは解放フラグを決してクリアせず，1ラウンドにつき1回だけ
反転させる．各スレッドは，入口で反転させるスレッド固有の変数
\texttt{localSense} を持つ．こうして \texttt{localSense} は，現在のラウンドが
解放されたときにフラグが取る値になる．最後に到着したスレッドがその値を
フラグに書き込む．カウンタ，フラグ，ロックは前の実装と同じ共有変数である．
\needspace{12\baselineskip}
\begin{lstlisting}[language=C, numbers=none, morekeywords={shared,private}]
private int localSense = 0;     /* スレッドごとに1つ */
void barrier(void) {
  localSense = !localSense;     /* 自分のセンスを反転 */
  lock(&countLock);
  int arrived = ++count;
  unlock(&countLock);
  if (arrived == N) {           /* 最後: 開く */
    count = 0;
    release = localSense;
  } else
    while (release != localSense)
      ;                         /* 待つ */
}
\end{lstlisting}

フラグが初期化し直されることはないので，上で遅れたスレッドは，フラグを読む
のがどれほど遅くなっても自分のラウンドを解放した値をそこに見出し，バリアを
出る．他のスレッドがまだこのバリアで待っている間に次のバリアへ到達した
スレッドがいても害はない．そのスレッドはカウンタをインクリメントするだけで
あり，自身の \texttt{localSense} は今や反対の値を保持しているので，次の反転を
待つ．連続するラウンドは2つの値を交互に取るので，スレッドが次のラウンドの
解放を自分のものと取り違えることはありえない．フラグを反転させる代わりに
ラウンドに番号を付けるのも同等の一般化であり，センスはラウンド番号の偶奇に
あたる．\cref{fig:l20-barrier} では，スレッド~2 が最初のラウンドを完了した
ときに \texttt{release} が 1 になり，スレッド~0 が2回目のラウンドを完了した
ときに再び 0 になる．

\begin{keypoint}[章のまとめ]
  共有データを更新するスレッドは同期しなければならない．さもなければ
  タイミング次第で更新が失われる．ロックはクリティカルセクションに相互排除
  を与え，1つの変数へのアクセスはすべて同じロックを使わなければならない．
  通常のロードとストアで作った素朴なロックは正しくない．その検査と設定自体
  がクリティカルセクションを必要とするからである．Lamport のパン屋の
  アルゴリズムのようなソフトウェアによる解法は正しいが遅すぎるので，
  プロセッサはアトミック命令を提供する．アトミック交換とテストアンドセット
  は1つの命令でメモリを読み書きするが，これはロード/ストア型のパイプライン
  には収まりが悪い．LL/SC は通常のメモリアクセス2回と，その間にその位置が
  書かれていればストアを失敗させるリンクレジスタとを用い，ロックなしで短い
  読み出し-変更-書き込みをアトミックにもする．スピンロックは，アトミックな
  書き込みの前に通常の読み出しでロック変数を検査すべきである．そうすれば
  待機中のコアはブロックをキャッシュ間で移動させずに共有状態に保つ．
  unlock は通常のストアである．バリア同期はスレッドの一群を互いに
  待ち合わせさせる．再利用されるバリアは，遅れたスレッドがまだ見ていない
  フラグをリセットするのではなく，各ラウンドを固有のフラグ値で，ラウンド
  ごとに反転させながら解放しなければならない．
\end{keypoint}

\end{document}
